% from ega4-18-fr.tex
\oldpage[IV-4]{212}

\begin{enumerate}
  \item[(i)] Considérons l'homomorphisme de degré $0$
  \[
    \psi_A:(\operatorname{gr}_{\mathfrak R}^*(A)\otimes_A(A/\mathfrak J))[T_1,\ldots,T_n]\longrightarrow\operatorname{gr}_\bullet'(A)=\operatorname{gr}_{\mathfrak L}^*(A)
  \]
  ($0$, 15.1.5.2). Le fait que cet homomorphisme soit \emph{surjectif} entraîne que la sous-$\operatorname{gr}_{\mathfrak L}^*(A)$-algèbre de $\operatorname{gr}_\bullet'(R)$ engendrée par $S'=\psi_R'(S)$ n'est autre que $\operatorname{Im}(\psi_R')$, d'où l'équivalence de $c)$ et $c')$.
  \item[(ii)] Puisque $e)$ est vérifiée, il en est de même de $c)$, et en vertu de (i), il reste à prouver la suffisance de la condition de l'énoncé. Or, comme $\psi_R'$ est bijectif, dire que $S'$ engendre $\operatorname{gr}_\bullet'(R)$ considéré comme $\operatorname{gr}_{\mathfrak L}^*(A)$-module équivaut à dire que $S$ engendre $(\operatorname{gr}_\bullet(R)\otimes_A(A/\mathfrak J))[T_1,\ldots,T_n]$ comme $\operatorname{gr}_{\mathfrak L}^*(A)$-module, et en vertu de la surjectivité de $\psi_A$, on peut dans cet énoncé remplacer l'anneau $\operatorname{gr}_{\mathfrak L}^*(A)$ par $(\operatorname{gr}_{\mathfrak R}^*(A)\otimes_A(A/\mathfrak J))[T_1,\ldots,T_n]$. Il revient évidemment au même de dire que $S$ engendre $\operatorname{gr}_\bullet(R)\otimes_A(A/\mathfrak J)$ en tant que module sur l'anneau $\operatorname{gr}_{\mathfrak R}^*(A)\otimes_A(A/\mathfrak J)$, ou aussi en tant que $\operatorname{gr}_{\mathfrak R}^*(A)$-module. Or, si $T$ est le sous-$\operatorname{gr}_{\mathfrak R}^*(A)$-module gradué de $\operatorname{gr}_\bullet(R)$ engendré par $S$, et si l'on pose $T'=\operatorname{gr}_\bullet(R)/T$, on a par hypothèse $T'/\mathfrak JT'=0$, ou $\mathfrak JT'=T'$; mais $T'$ est un module gradué, et chaque $T_p'$ est un quotient de $\operatorname{gr}_p(R)$, donc séparé pour la topologie $\mathfrak J$-préadique; l'hypothèse $\mathfrak JT_p'=T_p'$ entraîne donc $T_p'=0$ pour tout entier $p$, donc $T'=0$, ce qui achève de prouver le corollaire.
\end{enumerate}

\subsection{Critère de platitude normale de Hironaka.}
\EGAReaderTarget{ega4.c04.subsec.19-7}{subsection:IV.19.7}
\label{IV.19.7}

\begin{theorem}[19.7.1]
\label{IV.19.7.1}
\textup{(Hironaka).} Soient $A$ un anneau, $\mathfrak R$ un idéal de $A$, $\mathbf f=(f_1,\ldots,f_n)$ une suite d'éléments de $A$ qui est $(A/\mathfrak R)$-régulière, $\mathfrak J=f_1A+\cdots+f_nA$ l'idéal engendré par $\mathbf f$, $\mathfrak L=\mathfrak J+\mathfrak R$, $M$ un $A$-module.

On a dans ce cas
\[
  (\operatorname{gr}_{\mathfrak R}^*(M))\otimes_A(A/\mathfrak J)
  =(\operatorname{gr}_{\mathfrak R}^*(M))\otimes_A(A/\mathfrak L)
  =(\operatorname{gr}_{\mathfrak R}^*(M))\otimes_{A/\mathfrak R}(A/\mathfrak L),
\]
car
\[
  (\mathfrak R^nM/\mathfrak R^{n+1}M)\otimes_A(A/\mathfrak J)
  =\mathfrak R^nM/(\mathfrak R^{n+1}+\mathfrak J\mathfrak R^n)M
  =\mathfrak R^nM/(\mathfrak R^{n+1}+\mathfrak L\mathfrak R^n)M,
\]
et la dernière égalité résulte de Bourbaki, \emph{Alg.}, chap. II, 3\textsuperscript{e} éd., § 3, n\textsuperscript{o} 7, cor. 2 de la prop. 6.

Considérons les conditions suivantes:
\begin{enumerate}
  \item[a)] $\operatorname{gr}_{\mathfrak R}^*(M)$ est un $(A/\mathfrak R)$-module plat.
  \item[b)] $\operatorname{gr}_{\mathfrak L}^*(M)$ est un $(A/\mathfrak L)$-module plat et l'homomorphisme canonique ($0$, 15.1.5.2)
  \[
    \psi_M:((\operatorname{gr}_{\mathfrak R}^*(M))\otimes_A(A/\mathfrak L))[T_1,\ldots,T_n]\longrightarrow\operatorname{gr}_{\mathfrak L}^*(M)
  \]
  est bijectif.
  \item[c)] $(\operatorname{gr}_{\mathfrak R}^*(M))\otimes_{A/\mathfrak R}(A/\mathfrak L)$ est un $(A/\mathfrak L)$-module plat et la suite $\mathbf f$ est $(\operatorname{gr}_{\mathfrak R}^*(M))$-régulière.
  \item[d)] $\operatorname{gr}_{\mathfrak L}^*(M)$ est un $(A/\mathfrak L)$-module plat, et pour tout $\mathfrak p\in\operatorname{Ass}_{A/\mathfrak R}(A/\mathfrak L)$, $(\operatorname{gr}_{\mathfrak R}^*(M))_{\mathfrak p}$ est un $(A/\mathfrak R)_{\mathfrak p}$-module plat.
\end{enumerate}
On a alors les implications
\[
\xymatrix@C=3em@R=1.6em{
  a) \ar@{=>}[r] \ar@{=>}[dr] & c) \ar@{=>}[r] & b) \\
  & d)
}
\]
Lorsque tout quotient d'un sous-module d'un $\operatorname{gr}_{\mathfrak R}^p(M)$ ($p\geq0$) est idéalement séparé pour $\mathfrak J$ (Bourbaki, \emph{Alg. comm.}, chap. III, § 5, n\textsuperscript{o} 1), les conditions $a)$, $b)$ et $c)$ sont équivalentes.

